\section{Benchmark Construction}
\label{sec:benchmark}

\subsection{Dataset: EDA-Scripting-Bench}

\textbf{EDA-Scripting-Bench} is the instrument for both of this paper's comparisons: within one toolchain, the effect of each grounding component, and across toolchains, the effect of language structure. It covers PyAether, SKILL, and Tcl on four commercial runtimes, \textbf{449 tasks} in four slices (Table~\ref{tab:benchmark-composition}), all specified by one schema---a target language and runtime, a natural-language prompt, a function signature, and language-appropriate assertions---so that a single harness (Section~\ref{sec:agent-workflow}) scores all four platforms uniformly. Two properties carry the validity of these comparisons: scoring is execution-based, so a task passes only when the generated code acts correctly in the live runtime; and the cross-language slices are difficulty-controlled by construction (Section~\ref{sec:benchmark-slices}), so a score gap between slices reflects the toolchain rather than the tasks.
% EDA-Scripting-Bench 是本文两组比较的测量仪器：工具链内部——各grounding组件的作用；跨工具链——语言结构的影响。
% 覆盖PyAether/SKILL/Tcl三语言、四种商业运行时、四切片共449个任务（表~\ref{tab:benchmark-composition}），
% 全部由单一schema规定——目标语言与运行时、自然语言提示、函数签名、语言适配断言——使单一框架
% （Section~\ref{sec:agent-workflow}）统一评分。两点性质支撑比较的效度：评分基于执行，代码必须在实时
% 运行时中正确动作才算通过；跨语言切片在构造上控制了难度（Section~\ref{sec:benchmark-slices}），
% 切片间分差反映工具链而非任务本身。

\begin{table}[htbp]
\centering
\small
\caption{Composition of EDA-Scripting-Bench}
\label{tab:benchmark-composition}
\begin{tabular}{@{}llr@{}}
\toprule
\multicolumn{3}{@{}l}{\emph{Slices and runtimes}}\\
\textbf{Slice} & \textbf{Runtime} & \textbf{Tasks} \\
\midrule
EDA-Eval-PyAether & Empyrean Aether & 158 \\
EDA-Eval-Skill & Cadence Virtuoso & 157 \\
EDA-Eval-WaveViewTcl & Synopsys Custom WaveView & 77 \\
EDA-Eval-InnovusTcl & Cadence Innovus & 57 \\
\midrule
\multicolumn{2}{l}{Total} & \textbf{449} \\
\bottomrule
\end{tabular}
\end{table}
% 表：EDA-Scripting-Bench组成（切片 / 运行时 / 任务数）。PyAether 切片的场景分布已并入正文文字描述。

\subsection{The PyAether slice (EDA-Eval-PyAether)}

The reference slice, \textbf{EDA-Eval-PyAether} (158 tasks), anchors every task to a concrete PyAether API; EDA experts curated each assertion to be self-contained and executable in a clean sandbox without vendor-internal context. Tasks go beyond isolated single-API examples: they span layout, schematic, and design-management workflows---layout (64 tasks, 40.5\%) and schematic (47, 29.7\%) dominate the scenario mix, reflecting PyAether's role in physical design---compose several distinct APIs, and sometimes depend on a technology library, so passing reflects genuine command of the tool's behavior rather than surface-level retrieval.
% 参照切片EDA-Eval-PyAether（158任务）每个任务锚定一个具体PyAether API；EDA专家确保每条断言
% 自包含、可在无厂商内部上下文的干净沙盒中执行。任务超出孤立单API示例：覆盖layout/schematic/
% design-management工作流（layout 64任务/40.5%与schematic 47任务/29.7%主导场景，反映PyAether在物理设计中的核心角色），
% 组合多个不同API，部分依赖工艺库，故通过反映对工具行为的
% 真正掌握而非浅层检索。


\subsection{SKILL and Tcl slices}
\label{sec:benchmark-slices}

The remaining slices are built so that cross-language scores are comparable. \textbf{EDA-Eval-Skill} (157 SKILL tasks on Cadence Virtuoso) is \emph{function-aligned} with PyAether: each task re-implements, in Lisp style, the same entry point and functional specification as its PyAether counterpart (e.g., \texttt{create\_inverter\_schematic} builds the same transistor-level cell in both), holding task difficulty fixed; PyAether alone adds one pure-Python utility task with no EDA API. The \textbf{Tcl} slices are graded by difficulty and span two further toolchains: waveform measurement and signal querying on WaveView (\textbf{EDA-Eval-WaveViewTcl}, 77 tasks) and floorplanning, placement, and routing on Innovus (\textbf{EDA-Eval-InnovusTcl}, 57 tasks). Together these controls let the cross-language comparison (Section~\ref{sec:cross-language}) attribute score differences to language structure rather than task hardness.
% 其余切片为跨语言可比性而构建。EDA-Eval-Skill（157个Virtuoso上的SKILL任务）与PyAether
% 函数级对齐：每个SKILL任务以Lisp风格重实现对应PyAether任务的入口点与功能规范
% （如create_inverter_schematic在两者中构建同一晶体管级单元），固定任务难度；PyAether多出
% 1个无EDA API的纯Python工具任务。Tcl切片按难度分级，覆盖另外两种工具链：WaveView上的波形
% 测量与信号查询（EDA-Eval-WaveViewTcl，77任务）与Innovus上的floorplan/placement/routing
% （EDA-Eval-InnovusTcl，57任务）。这些控制使跨语言比较（Section~\ref{sec:cross-language}）
% 将分差归因于语言结构而非任务难度。

\subsection{Evaluation Protocol}

A task passes if the generated code executes in its target runtime sandbox without errors and all assertions hold; the primary metric is Pass@1 \cite{chen2021evaluating}:
% 若生成的代码在其目标运行时沙盒中无错误执行且所有断言成立，则该任务通过；主要指标为Pass@1 \cite{chen2021evaluating}：

\[
\text{Pass@1} = \frac{\text{\# tasks with all assertions passed}}{\text{Total tasks}} \times 100\%
\]
% \[
% \text{Pass@1} = \frac{\text{所有断言均通过的任务数}}{\text{总任务数}} \times 100\%
% \]